\section{Symmetry-adapted Fock encoding}

\subsection{Circular vibrational bases and experimental accessibility}
\label{app:circular_vibrations}

The circular basis is a unitary change of basis within a degenerate vibrational manifold and requires no modification of the molecule. Although the normal-coordinate angular momentum $L_Q=Q_1P_2-Q_2P_1$ generates rotations in the degenerate mode space, its relation to physical molecular angular momentum depends on the underlying atomic displacement patterns \cite{Herzberg1942,Amat1958}.

Circular vibrational excitation requires coherent driving of the two degenerate components with a relative phase of $\pm\pi/2$. For an infrared-active doublet, the interaction
\begin{equation}
H_{\mathrm{int}}
=-\sum_{j=1}^{2}
[\mathbf E(t)\cdot\mathbf d_j]Q_j,
\qquad
\mathbf d_j=
\left.\frac{\partial\boldsymbol{\mu}}{\partial Q_j}\right|_0,
\end{equation}
allows such excitation using resonant circularly polarized light when the transition-dipole derivatives have suitable orientations and strengths. Phase-controlled fields can accommodate more general couplings, while Raman excitation provides an alternative for Raman-active modes. Coherent vibrational excitation and time-resolved optical detection have been demonstrated experimentally \cite{Peschel2022,CoherentIR2024}. Circularly polarized molecular vibrations and their polarization-resolved pump--probe detection have been investigated theoretically \cite{Lin2025}.







\subsection{Proof of Corollary}~\ref{cor:shift}
\label{proof:shift}
Since \(\Gvib\) is one-dimensional, \(V_{\Grot}\otimes V_{\Gvib}\) is an irreducible carrier of species \(\Grot\otimes\Gvib\). Regrouping and exchanging tensor factors gives the natural \(\grp{G}\)-equivariant isomorphism
\[
V_{\Grot}\otimes V_{\Gspin}
\otimes V_{\Gvib}
\cong
V_{\Grot\otimes\Gvib}
\otimes V_{\Gspin}\otimes V_{\Gamma_0},
\]
with \(\Gamma_0\) the trivial irrep of \(\grp{G}\).
It therefore identifies the corresponding \(\Gmol\)-coupling spaces, proving equal multiplicities and simultaneous admissibility. Matching $V_{\Grot}\otimes V_{\Gvib}$ to the standard carrier of $\Grot\otimes\Gvib$ takes a fixed intertwiner, for example $\diag(1,-1)$ for $\irr{e}\otimes\irr{a}_{2}$ in $\grp{C_{3v}}$. After this choice the identification is fixed up to a phase. The nuclear-spin factor is unchanged, so its statistical weight is unchanged.

For Abelian \(\grp{G}\), all carriers are one-dimensional. Admissibility is therefore equivalent to \(
\Grot\otimes\Gspin\otimes\Gvib
=\Gmol,
\)
which has multiplicity one and yields \(\Grot
=\Gmol\otimes\Gspin^*
\otimes\Gvib^*\).
\hfill$\square$


\subsection{Vibrational states of calcium monomethoxide}
\label{app:caoch3}

$\mathrm{CaOCH_3}$ has six nuclei and 12 vibrational degrees of freedom, transforming as $4\irr{a}_{1}\oplus4\irr{e}$ under $\grp{C_{3v}}$, in the frame of Table~\ref{tab:molecular_examples}. With four doubly degenerate sets the vibrational Fock space is richer than that of phosphine, but a single rule assigns all of its carriers. On a product of circular occupation states of the $\irr{e}$ modes, $\bigotimes_s\occ{n_{s-}n_{s+}}$, the threefold rotation acts as $\rootunity{3}^{-\ell}$ with total vibrational angular momentum $\ell=\sum_s(n_{s+}-n_{s-})$, and $\elt{\sigma_{v}}(0)$ exchanges $n_{s+}\leftrightarrow n_{s-}$ in every mode simultaneously. Consequently a product state with $\ell\equiv2\pmod 3$ is a $\irket{\irr{e}}{1}$ carrier with partner $\irket{\irr{e}}{2}=\elt{\sigma_{v}}(0)\irket{\irr{e}}{1}$, a state with $\ell\equiv0$ and not invariant under $\elt{\sigma_{v}}(0)$ generates one $\irr{a}_{1}$ and one $\irr{a}_{2}$ state as $\tfrac{1}{\sqrt{2}}\,(1\pm\elt{\sigma_{v}}(0))$, and a $\elt{\sigma_{v}}(0)$-invariant state with $\ell=0$ is $\irr{a}_{1}$.

For the Ca--O--C bend $\nu_8$ this gives
\begin{align}
v_8=1:&\quad \irket{\irr{e}}{1}=\occ{10},\ \irket{\irr{e}}{2}=\occ{01};\nonumber\\
v_8=2:&\quad \ket{\irr{a}_{1}}=\occ{11},\ \irket{\irr{e}}{1}=\occ{02},\ \irket{\irr{e}}{2}=\occ{20};\nonumber\\
v_8=3:&\quad \ket{\irr{a}_{1}},\ket{\irr{a}_{2}}=\tfrac{1}{\sqrt{2}}\,(\occ{03}\pm\occ{30}),\ \irket{\irr{e}}{1}=\occ{21},\ \irket{\irr{e}}{2}=\occ{12},
\label{eq:caoch3_bend}
\end{align}
in the notation $\occ{n_{8-}n_{8+}}$. The combination band $\nu_7+\nu_8$ gives $\irr{e}\otimes\irr{e}=\irr{a}_{1}\oplus\irr{a}_{2}\oplus\irr{e}$ with $\ket{\irr{a}_{1}},\ket{\irr{a}_{2}}=\tfrac{1}{\sqrt{2}}\,(\ket{\ell_7{=}{+1},\ell_8{=}{-1}}\pm\ket{\ell_7{=}{-1},\ell_8{=}{+1}})$ and $\irket{\irr{e}}{1}=\ket{\ell_7{=}{+1},\ell_8{=}{+1}}$, $\irket{\irr{e}}{2}=\ket{\ell_7{=}{-1},\ell_8{=}{-1}}$. Every vibrational state of the molecule is thereby assigned a carrier, and the admissible species and coupling tensors are those of Table~\ref{tab:ammonia_species}. The $\irr{a}_{2}$ state at $v_8=3$, at $432\ \mathrm{cm^{-1}}$ in the harmonic approximation, is the lowest vibrational level that exchanges the $\irr{a}_{1}$ and $\irr{a}_{2}$ rotational partners of ortho spin (Corollary~\ref{cor:shift}).



\section{Molecular Fourier Transform}


\subsection{Loop monodromy and molecular symmetry}
\label{app:monodromy}

A closed loop in the molecular orientation space \(\Othree/\grp{G}\), when lifted from an orientation \(\orient\), may terminate at \(\orient\gel\), where \(\gel\in \grp{G}\). Because the lifted path remains within a single connected component of \(\Othree\), the corresponding group element must belong to the proper-rotation subgroup \(\grp{G}\cap\SOthree\). The associated monodromy is represented by \(D^{\Grot}(\gel)\). Consequently, non-Abelian loop monodromy requires a noncommutative image of \(\grp{G}\cap\SOthree\) under \(D^{\Grot}\). A multidimensional representation alone is insufficient. This discussion concerns symmetry-induced flat holonomy. Berry and dynamical connections, when present~\cite{Berry1984, Mead1992}, must be treated separately.







\subsection{Proof of Theorem~\ref{thm:fourier}} 
\label{proof:fourier}

Expand a sharp orientation in Wigner functions, $$\ket{\orient}=\sum_{J,\parity,m,k}\sqrt{\frac{2J+1}{16\pi^2}}\,\irrJp{J}{\parity}_{m,k}(\orient)^*\ket{{}^{J,\parity}_{\,m},k},$$ insert it into Eq.~\eqref{eq:coset_states} with $\orient=\crep\gel$, and pass to the $\grp{G}$-adapted basis $\cgket{\Gamma}{\nu}{{}^{J,\parity}_{\,m}\mult}$ (partner $\nu$, copy $\mult$) in which $\cev X_{\gel}$ acts as $D^{\Gamma}_{\nu',\nu}(\gel)\delta_{\mult\mult'}$. Splitting $\irrJp{J}{\parity}(\crep\gel)=\irrJp{J}{\parity}(\crep)\irrJp{J}{\parity}(\gel)$ and applying the Schur orthogonality relation $$\sum_{\gel} D^{\Grot}_{\lambda,\mu}(\gel)D^{\Gamma}_{\nu',\nu}(\gel)^{*}=\frac{|\grp{G}|}{d_{\Grot}}\delta_{\Gamma\Grot}\delta_{\lambda\nu'}\delta_{\mu\nu}$$ collapses the group sum to
\begin{equation}
\ket{\crep;\Grot^{\lambda\mu}}=\sum_{m^{(\mathrm r)}_{\Grot}(J,\parity)>0}\sum_{m,\mult}\sqrt{\frac{(2J+1)|\grp{G}|}{16\pi^2 d_{\Grot}}}\,\irrJp{J}{\parity}_{m;\lambda\mult}(\crep)^*\,\cgket{\Grot}{\mu}{{}^{J,\parity}_{\,m}\mult}.
\end{equation}
The prefactor is fixed by $\sqrt{\frac{2J+1}{16\pi^2}}\cdot\sqrt{\frac{d_{\Grot}}{|\grp{G}|}}\cdot\frac{|\grp{G}|}{d_{\Grot}}$, which is Eq.~\eqref{eq:induced-harmonic}. It reduces to $\sqrt{\frac{(2J+1)|\grp{G}|}{8\pi^2 d}}$ of Ref.~\cite{albert2025quantumtheorymolecularorientations} on replacing $\grp{O(3)}$ by $\grp{SO(3)}$. Contracting the index $\mu$ with the coupling tensor, which acts on neither $(\crep,\lambda)$ nor $(m,\mult)$, gives Eq.~\eqref{eq:ft_forward}. For orthogonality, $\sum_\lambda|\irrJp{J}{\parity}_{m;\lambda\mult}(\crep)|^2$ is invariant under $\crep\to{\crep}\gel$, so its integral over the quotient is $|\grp{G}|^{-1}$ times its integral over $\grp{O(3)}$, which by the Peter--Weyl orthogonality of $\grp{O(3)}$ is $d_{\Grot}\cdot16\pi^2/(2J+1)$. The product with the squared prefactor is one, and the off-diagonal terms vanish by the same argument. Completeness follows from completeness of the Wigner functions on $\grp{O(3)}$ restricted to the $\Grot$-covariant subspace. \hfill$\square$





\subsection{Symmetry operation on different state spaces}
\label{subsec:tour}

Before the full species tables, it is useful to watch what a single symmetry operation does to a concrete state and how the answer changes as sectors are added one at a time. We use phosphine and follow the two generator operations of $\grp{C_{3v}}$: the threefold rotation $\elt{C_{3}}=\rotax{z}(2\pi/3)$, which carries proton $j$ (at azimuth $2\pi(j-1)/3$) to proton $j+1$, and the reflection $\elt{\sigma_{v}}(0)=\elt{i}\,\rotax{y}(\pi)$ in the $xz$ plane, which fixes proton $1$ and exchanges protons $2$ and $3$. We take four state spaces in turn, all at the lowest rotational and vibrational excitation that shows the effect: (i) vibration alone; (ii) rotation and vibration, the rovibrational problem of conventional symmetric-top spectroscopy; (iii) rotation and nuclear spin, the two sectors of Ref.~\cite{albert2025quantumtheorymolecularorientations}; and (iv) all three. Table~\ref{tab:tour} summarizes the four steps. Throughout, $\rootunity{3}=\expe^{2\pi\iunit/3}$.

\emph{(i) Vibration alone.} One quantum in the degenerate bend $\nu_4$ spans a two-dimensional space. In the Cartesian basis, Eq.~\eqref{eq:creation_transform} gives the matrices of the normal-coordinate transformation itself, 
$U(\elt{C_{3}})
=
\frac12\left(\begin{smallmatrix}-1&-\sqrt3\\ \sqrt3&-1\end{smallmatrix}\right),\quad
U(\elt{\sigma_{v}}(0))=\left(\begin{smallmatrix}1&0\\0&-1\end{smallmatrix}\right).$
The rotation does not multiply the two components by a phase. Instead, it turns one displacement pattern into a superposition of both, so no single Cartesian component of $\nu_4$ has a definite symmetry. In the circular basis $\ket{\ell_4{=}{\mp1}}$ the same operations become
\begin{equation}
U(\elt{C_{3}})=\left(\begin{smallmatrix}\rootunity{3}&0\\0&\rootunity{3}^{2}\end{smallmatrix}\right),\qquad
U(\elt{\sigma_{v}}(0))=\left(\begin{smallmatrix}0&1\\1&0\end{smallmatrix}\right),
\label{eq:tour_vib_circ}
\end{equation}
so $\elt{C_{3}}$ multiplies a state of vibrational angular momentum $\ell$ by $\rootunity{3}^{-\ell}$ and $\elt{\sigma_{v}}(0)$ reverses the sense of circulation. The characters $(2,-1,0)$ on $\{\cls{E},\cls{2C_{3}},\cls{3\sigma_{v}}\}$ identify the carrier $\irr{e}$. 

Two quanta span $\{\ket{\ell_4{=}{+2}},\ket{\ell_4{=}0},\ket{\ell_4{=}{-2}}\}$, on which $U(\elt{C_{3}})=\diag(\rootunity{3},1,\rootunity{3}^{2})$ and $\elt{\sigma_{v}}(0)$ exchanges $\ell_4=\pm2$ and fixes $\ell_4=0$. The characters $(3,0,1)$ give $\operatorname{Sym}^{2}\irr{e}=\irr{a}_{1}\oplus \irr{e}$, the $\ell_4=0$ state being totally symmetric. For an experiment this is the familiar infrared selection rule: the dipole transforms as $\irr{a}_{1}\,(\mu_z)\oplus \irr{e}\,(\mu_x,\mu_y)$, so the $\nu_4$ fundamental is a perpendicular band, while $2\nu_4$ has a parallel ($\ell_4=0$) and a perpendicular ($\ell_4=\pm2$) component. What vibration alone cannot supply is a statistical weight: without nuclear spin there is nothing for the spin-statistics condition~\eqref{eq:gamma_mol} to act on.

\emph{(ii) Rotation and vibration.} We add the rotation at $J=1$, $\parity=+1$. By Eq.~\eqref{eq:rot_action} the molecule-frame operations act on the projection $k$ along the symmetry axis as
\begin{equation}
\elt{C_{3}}\ket{{}^{1,+}_{\,m},k}=\rootunity{3}^{k}\ket{{}^{1,+}_{\,m},k},\qquad
\elt{\sigma_{v}}(0)\ket{{}^{1,+}_{\,m},k}=(-1)^{1+k}\ket{{}^{1,+}_{\,m},-k},
\label{eq:tour_rot}
\end{equation}
while the laboratory projection $m$ is untouched. On a product state the two phases combine: $$\elt{C_{3}}\ket{k}\ket{\ell_4}=\rootunity{3}^{k-\ell_4}\ket{k}\ket{\ell_4},$$ and $\elt{\sigma_{v}}(0)$ sends $(k,\ell_4)\to(-k,-\ell_4)$, here with sign $+1$ for $|k|=1$. The four states with $|k|=|\ell_4|=1$ therefore split into two pairs. The pair $(k,\ell_4)=(1,1),(-1,-1)$ has $k-\ell_4=0$, is invariant under $\elt{C_{3}}$, and is exchanged by $\elt{\sigma_{v}}(0)$, so it gives
\begin{equation}
\begin{aligned}
    \tfrac{1}{\sqrt{2}}\,\bigl(\ket{1,{+1}}+\ket{{-1},{-1}}\bigr)\in \irr{a}_{1},\\
    \tfrac{1}{\sqrt{2}}\,\bigl(\ket{1,{+1}}-\ket{{-1},{-1}}\bigr)\in \irr{a}_{2},
\end{aligned}
\label{eq:tour_rovib}
\end{equation}
whereas the pair $(1,-1),(-1,1)$ acquires $\rootunity{3}^{2}$ and $\rootunity{3}$ and forms an $\irr{e}$ carrier. The $k=0$ state ($\irr{a}_{2}$) times $\nu_4$ ($\irr{e}$) gives a second $\irr{e}$. The six-dimensional space has characters $(6,0,0)$ and decomposes as $\irr{a}_{1}\oplus \irr{a}_{2}\oplus2\irr{e}$. This is the rule $k-\ell\equiv0\pmod 3$ for $\irr{a}$-type rovibrational levels, and Eq.~\eqref{eq:tour_rovib} shows where the $\ell$-type doubling of the $k\ell>0$ levels comes from: those levels form an $\irr{a}_{1},\irr{a}_{2}$ pair that a symmetry-allowed term can split, whereas the $k\ell<0$ levels form a degenerate $\irr{e}$ pair that no symmetric term can split. The perpendicular-band selection rule $\Delta k=\Delta\ell=\pm1$ with equal signs is the statement that the dipole preserves $k-\ell\bmod3$. Again, however, this space carries no spin, so it cannot tell which of these levels exist.

\emph{(iii) Rotation and nuclear spin.} Return to the vibrational ground state and add the three proton spins, the setting of Ref.~\cite{albert2025quantumtheorymolecularorientations}. Now each operation also permutes nuclei: $\elt{C_{3}}$ cycles the spins, $\ket{s_1s_2s_3}\to\ket{s_3s_1s_2}$, and $\elt{\sigma_{v}}(0)$ transposes the spins of protons $2$ and $3$. Because the protons are fermions, a physical state must be multiplied by $\chi_{\irr{a}_{2}}=(+1,+1,-1)$, i.e.\ it must be invariant under $\elt{C_{3}}$ and change sign under $\elt{\sigma_{v}}(0)$. Three trial states show how the rule operates at $J=1$, $\parity=+1$.
\begin{itemize}
\item $\ket{{}^{1,+}_{\,m},0}\otimes\ket{\uparrow\uparrow\uparrow}$: $\elt{C_{3}}$ gives $1\cdot1$ and $\elt{\sigma_{v}}(0)$ gives $(-1)\cdot(+1)=-1$. It is physical: an ortho state, one of four, since all four $I_{\mathrm{tot}}=\tfrac32$ spin states behave alike.
\item $\ket{{}^{1,+}_{\,m},0}\otimes\cgket{\irr{e}}{1}{\tfrac12\tfrac12}$: $\elt{C_{3}}$ gives $\rootunity{3}\neq1$. It is not physical, and no linear combination with other $k=0$ states can make it so.
\item $\tfrac{1}{\sqrt{2}}\,\bigl(\ket{{}^{1,+}_{\,m},1}\cgket{\irr{e}}{2}{\tfrac12M}-\ket{{}^{1,+}_{\,m},{-1}}\cgket{\irr{e}}{1}{\tfrac12M}\bigr)$: $\elt{C_{3}}$ gives $\rootunity{3}\cdot\rootunity{3}^{2}=1$, and $\elt{\sigma_{v}}(0)$ exchanges the two terms and so gives $-1$. It is physical: a para state, one for each $M=\pm\tfrac12$.
\end{itemize}
The $24$-dimensional space has characters $(24,0,-4)$ and contains $\irr{a}_{2}$ six times: four ortho states at $k=0$ and two para states at $|k|=1$, per laboratory projection. At $J=0$ the same count gives no physical state for $\parity=+1$ and four for $\parity=-1$. These are the observables of this step: nuclear-spin statistical weights ($4$ for $K=|k|=0,3,6,\ldots$, $2$ for the other $K$, before the $^{31}$P factor of $2$) and missing levels (only one parity component of each $K=0$ level exists).

\emph{(iv) Rotation, nuclear spin, and vibration.} Finally, excite one $\nu_4$ quantum with all three sectors present. The perrotation is the tensor product of the three actions above, $\cev X_{\gel}\otimes\mathrm{Perm}(\gel)\otimes U_{\mathrm{v}}(\gel)$, so on a product state the phases of all three factors multiply. The $48$-dimensional space at $J=1$, $\parity=+1$ has characters $(48,0,0)$ and contains $\irr{a}_{2}$ eight times. The eight physical states are:
\begin{itemize}
\item four \emph{ortho} states, $\tfrac{1}{\sqrt{2}}\,\bigl(\ket{1,{+1}}-\ket{{-1},{-1}}\bigr)\otimes\ket{\tfrac32,M_I}$: the $\irr{a}_{2}$ rovibrational combination of Eq.~\eqref{eq:tour_rovib}, species $(\irr{e},\irr{a}_{1},\irr{e})$;
\item two \emph{para} states of species $(\irr{a}_{2},\irr{e},\irr{e})$, $\ket{{}^{1,+}_{\,m},0}\otimes\tfrac{1}{\sqrt{2}}\,\bigl(\cgket{\irr{e}}{1}{\tfrac12M}\ket{\ell_4{=}{+1}}+\cgket{\irr{e}}{2}{\tfrac12M}\ket{\ell_4{=}{-1}}\bigr)$;
\item two \emph{para} states of species $(\irr{e},\irr{e},\irr{e})$, the GHZ states $\catT{1}$ of Eq.~\eqref{eq:nh3_ghz_momentum}.
\end{itemize}
Checking the last one directly: $\elt{C_{3}}$ multiplies $\ket{k{=}1}\irket{\irr{e}}{1}_{\mathrm{n}}\ket{\ell_4{=}{-1}}$ by $\rootunity{3}\cdot\rootunity{3}\cdot\rootunity{3}=1$, and $\elt{\sigma_{v}}(0)$ exchanges it with $\ket{k{=}{-1}}\irket{\irr{e}}{2}_{\mathrm{n}}\ket{\ell_4{=}{+1}}$, so the antisymmetric combination changes sign, as $\Gmol=\irr{a}_{2}$ requires. 

Step (iv) also changes what an experiment sees. In the vibrational ground state every $|k|=1$ level is para. With one quantum in $\nu_4$, the $|k|=1$ levels divide by the relative sign of $k$ and $\ell_4$: the $k\ell>0$ level becomes ortho, with only its $\irr{a}_{2}$ $\ell$-doubling component present (its $\irr{a}_{1}$ partner has no spin partner and is missing~\cite{Oka1973}), and the $k\ell<0$ level remains para. Counting only the proton factor, the statistical weights of the subbands of a perpendicular band therefore follow $k-\ell\bmod3$ rather than $k\bmod3$: weight $4$ when $k-\ell\equiv0$ and weight $2$ otherwise. This is the familiar strong--weak--weak alternation of $\Delta K=\pm1$ subbands~\cite{Hougen1962}, obtained here from the same projection that produces the entangled states. For a rigid $\grp{C_{3v}}$ top such as $\mathrm{PH_3}$ or $\mathrm{CaOCH_3}$ (Sec.~\ref{sec:experiment}) the counting is exactly as stated, while in ammonia each level with $k\neq0$ is in addition split into an inversion doublet.

\begin{table*}[t]
\caption{The four steps of Sec.~\ref{subsec:tour} for phosphine}, at $J=1$, $\parity=+1$ where rotation is present, per laboratory projection $m$. Characters are on $\{\cls{E},\cls{2C_{3}},\cls{3\sigma_{v}}\}$. ``Physical'' is the multiplicity of $\Gmol=\irr{a}_{2}$, which can be imposed only once nuclear spin is present.
\label{tab:tour}
\centering
\small
\begin{tabular}{lcccll}
\toprule
State space & dim & $\chi$ & decomposition & physical ($\irr{a}_{2}$) & what it predicts\\
\midrule
(i) $\nu_4$ & 2 & $(2,-1,0)$ & $\irr{e}$ & --- & perpendicular IR band\\
\phantom{(i)} $2\nu_4$ & 3 & $(3,0,1)$ & $\irr{a}_{1}\oplus \irr{e}$ & --- & parallel and perpendicular components\\
(ii) rotation $\otimes$ $\nu_4$ & 6 & $(6,0,0)$ & $\irr{a}_{1}\oplus \irr{a}_{2}\oplus2\irr{e}$ & --- & $k-\ell$ rule, $\ell$-type doubling\\
(iii) rotation $\otimes$ spin & 24 & $(24,0,-4)$ & $2\irr{a}_{1}\oplus6\irr{a}_{2}\oplus8\irr{e}$ & 6 (4 ortho, 2 para) & statistical weights, missing levels\\
(iv) rotation $\otimes$ spin $\otimes$ $\nu_4$ & 48 & $(48,0,0)$ & $8\irr{a}_{1}\oplus8\irr{a}_{2}\oplus16\irr{e}$ & 8 (4 ortho, 4 para) & weights set by $k-\ell$; GHZ state $\catT{1}$\\
\bottomrule
\end{tabular}
\end{table*}






\section{Experimental preparation of SALCs in calcium monomethoxide}
\label{app:caoch3_prep}

This section of appendix details the protocol in Sec.~\ref{sec:experiment} for preparing coherent superpositions associated with molecular symmetry-adapted linear combinations (SALCs) in calcium monomethoxide. Within the stated approximations, we derive the level energies and transition frequencies and the terms in the Hamiltonian (Appendix~\ref{app:caoch3_ham}), dipole matrix elements (Appendix~\ref{app:caoch3_dipoles}), and pulse sequence and experimental requirements (Appendix~\ref{app:caoch3_pulses}).

\subsection{Hamiltonian of calcium monomethoxide}
\label{app:caoch3_ham}

We calculate every level energy and transition wavenumber in this appendix from Eq.~\eqref{eq:caoch3_hamiltonian}~\cite{JohnstonDennison1935,Herzberg1945,Nielsen1951}, the same finite set of terms that the symmetry analysis in Sec.~\ref{sec:experiment} uses. Its constants combine experimental data and ab initio calculations, and the upper block of Table~\ref{tab:caoch3_constants} lists the values we use. The terms missing from Eq.~\eqref{eq:caoch3_hamiltonian} are not added to these energies. Instead, we estimate the size of each missing term and state what it requires of the pulses. The lower block of Table~\ref{tab:caoch3_constants} lists the constants used for these estimates.

\emph{Terms of Eq.~\eqref{eq:caoch3_hamiltonian}.} The $\nu_8$ fundamental and rotational constants $A$ and $B$ are measured~\cite{Jacox2003,Kozyryev2019,Namiki1998,Crozet2002}. We estimate the remaining parameters using PySCF~\cite{Sun2020PySCF}. Coupled-cluster normal modes give $\zeta_8=0.926$, while second-order vibrational perturbation theory gives $x_{8,8}=-0.53\ \mathrm{cm^{-1}}$ and the anharmonic contribution $g^{\mathrm{anh}}_{8,8}=0.63\ \mathrm{cm^{-1}}$~\cite{Nielsen1951,Watson1968}. The anharmonic estimates vary by up to $0.11\ \mathrm{cm^{-1}}$ with finite-difference step size and density functional. We include the bend-angular-momentum contribution through $g_{8,8}=A\zeta_8^2+g^{\mathrm{anh}}_{8,8}$~\cite{Watson1968} and choose $\tilde\omega_8=\nu_8-g_{8,8}$, so that Eq.~\eqref{eq:caoch3_hamiltonian} reproduces the measured fundamental. We use its central value, $\nu_8=144\ \mathrm{cm^{-1}}$. Its uncertainty of $\pm5\ \mathrm{cm^{-1}}$ shifts all seven transitions of the tree together and leaves their spacings unchanged.

\emph{Terms beyond Eq.~\eqref{eq:caoch3_hamiltonian}.} Two terms omitted in Sec.~\ref{sec:experiment} have measured ground-state constants. Neither enters the energies and wavenumbers quoted in this appendix, and we use each only as follows. Centrifugal distortion adds $-D_NN^2(N+1)^2-D_{NK}N(N+1)k^2-D_Kk^4$ to Eq.~\eqref{eq:caoch3_hamiltonian}~\cite{Namiki1998,Crozet2002}. Within the preparation tree, it would shift the energies relative to $G$ and the transition wavenumbers by at most $0.013\ \mathrm{cm^{-1}}$. This is about a tenth of the spread in the ab initio constants above, so adding it would not make the predicted lines more accurate, and we leave it out. Spin--rotation and proton hyperfine interactions~\cite{Namiki1998} split each rovibrational level into $J$ and $F$ components and shift the tree levels by at most $25$~MHz ($0.0008\ \mathrm{cm^{-1}}$). We omit these shifts from the energies as well. Each pulse must still address a single component, so we use this spin structure to set the Rabi frequency in Appendix~\ref{app:caoch3_pulses}.

The remaining terms have no measured constants for the bend levels, and we neglect them. We use ground-state rotational, distortion, and spin constants for all bend levels, neglecting their $v$ dependence. We also neglect $\ell$-type doubling ($\Delta k=\Delta\ell=\pm2$), whose $\nu_8$ constant is unmeasured~\cite{Nielsen1951}. In polar bend coordinates $(\rho,\vartheta)$, a cubic potential term $\propto\rho^3\cos3\vartheta$ mixes $\ell$ values differing by three. Its perturbative energy shifts, and those from cubic couplings to other modes, enter the constants of Eq.~\eqref{eq:caoch3_hamiltonian}, but we neglect the associated state mixing. All terms beyond Eq.~\eqref{eq:caoch3_hamiltonian} preserve molecular point group symmetry and can mix only levels of the same point group irreps~\cite{BunkerJensen2006}. Thus the species triples $(\Grot,\Gspin,\Gvib)$ in Table~\ref{tab:caoch3_levels} become approximate labels, while every tree level retains $\Gmol=\irr{a}_{2}$.

\begin{table}[h]
\centering
\small
\caption{Constants of the $\widetilde X\,{}^2A_1$ ground state, grouped by use. The upper block holds the constants of Eq.~\eqref{eq:caoch3_hamiltonian}; every level energy and transition wavenumber of the preparation tree is computed from these alone. The last two entries of this block are derived from the others and rounded here, and the calculations use the unrounded values. The lower block holds the constants of terms beyond Eq.~\eqref{eq:caoch3_hamiltonian}, which enter no tabulated energy: the distortion constants give only the size of the omitted shifts, and the spin constants set only the Rabi frequency $\Omega$ in Appendix~\ref{app:caoch3_pulses}. ``This work'' marks ab initio estimates; the other entries are literature values.}
\label{tab:caoch3_constants}
\begin{tabular}{lll}
\toprule
Quantity & Value & Source\\
\midrule
\multicolumn{3}{l}{\emph{Eq.~\eqref{eq:caoch3_hamiltonian}, used for all energies and wavenumbers}}\\
$\nu_8$ fundamental & $144\pm5\ \mathrm{cm^{-1}}$ & Refs.~\cite{Jacox2003,Kozyryev2019}\\
$A$ & $5.448303\ \mathrm{cm^{-1}}$ & Ref.~\cite{Crozet2002}\\
$B$ & $0.11626491\ \mathrm{cm^{-1}}$ ($3485.534(2)$~MHz) & Refs.~\cite{Namiki1998,Crozet2002}\\
$\zeta_8$ & $0.926$ & EA-EOM-CCSD, this work\\
$g_{8,8}=A\zeta_8^2+g^{\mathrm{anh}}_{8,8}$ & $4.672+0.63=5.30\ \mathrm{cm^{-1}}$ & $g^{\mathrm{anh}}_{8,8}$: B3LYP, this work\\
$x_{8,8}$ & $-0.53\ \mathrm{cm^{-1}}$ & B3LYP, this work\\
$\tilde\omega_8=\nu_8-g_{8,8}$ & $138.70\ \mathrm{cm^{-1}}$ & from $\nu_8$ and $g_{8,8}$\\
\midrule
\multicolumn{3}{l}{\emph{Terms beyond Eq.~\eqref{eq:caoch3_hamiltonian}, used only for estimates}}\\
$D_N$, $D_{NK}$ & $2.6885\times10^{-8}$, $2.3580\times10^{-6}\ \mathrm{cm^{-1}}$ & Refs.~\cite{Namiki1998,Crozet2002}\\
$D_K$ & $7.03\times10^{-5}\ \mathrm{cm^{-1}}$ & Ref.~\cite{Crozet2002}\\
$\varepsilon_{bc}$ (spin--rotation) & $12.451(11)$~MHz & Ref.~\cite{Namiki1998}\\
$a_F$, $T_{aa}$ (${}^1$H hyperfine) & $-0.421(27)$, $1.070(45)$~MHz & Ref.~\cite{Namiki1998}\\
\bottomrule
\end{tabular}
\end{table}
The preparation tree follows $k=\ell-3$, with $N=3$ or $4$ (Table~\ref{tab:caoch3_levels}). For a transition $(v,\ell,N,k)\to(v+1,\ell+q,N',k+q)$ with $q=\pm1$, Eq.~\eqref{eq:caoch3_hamiltonian} gives
\begin{equation}
\begin{aligned}
\tilde\nu={}&\tilde\omega_8+2x_{8,8}v+g_{8,8}(2q\ell+1)\\
&+B\bigl[N'(N'+1)-N(N+1)\bigr]+(A-B)(2qk+1)\\
&-2A\zeta_8\bigl[q(k+\ell)+1\bigr].
\end{aligned}
\label{eq:caoch3_lines}
\end{equation}
Table~\ref{tab:caoch3_pulses} lists the resulting wavenumbers of the seven transitions.

P1, P2, and P4 are consecutive $q=+1$ transitions with $\Delta N=0$. From one to the next, $v$, $\ell$, and $k$ of the lower level each increase by one, so Eq.~\eqref{eq:caoch3_lines} places these lines at equal spacings of $2x_{8,8}+2g^{\mathrm{anh}}_{8,8}+2A(1-\zeta_8)^2-2B=0.027\ \mathrm{cm^{-1}}$, or $0.81$~GHz. The anharmonic part, $+0.20\ \mathrm{cm^{-1}}$, and the rotational and Coriolis part, $-0.17\ \mathrm{cm^{-1}}$, nearly cancel. The spacing is therefore smaller than the spread in the ab initio constants and sensitive to the neglected $v$ dependence of the constants, so these lines require further spectroscopic measurements.

Competing transitions that start or end on a populated level lie at least $0.73$~GHz from every tree line, with one exception. In Eq.~\eqref{eq:caoch3_lines}, a line with $\Delta N=0$ does not depend on $N$. P4 therefore coincides with its $N=4$ counterpart, the transition between the level with $v=2$, $\ell=2$, $k=-1$, $N=4$ and $L_{30}$. Terms beyond Eq.~\eqref{eq:caoch3_hamiltonian} lift this coincidence: centrifugal distortion alone separates the two lines by $0.57$~MHz, and the unmeasured $v$ dependence of $B$ adds to it. Because this separation is small and uncertain, the protocol does not rely on it: P4 must finish before P5 populates $L_{30}$ (Appendix~\ref{app:caoch3_pulses}).

For transition strengths, let $\ket{k}$ denote the rotational state at the specified $N$, and $\ket{\irr{a}_{1}}_{\mathrm n}$ the ortho nuclear-spin state with $I=\tfrac32$. The circular bend occupations $n_\pm$ give $v=n_++n_-$ and $\ell=n_+-n_-$. Suppressing spin-coupling labels, each tree state combines two orthogonal reflection partners:
\begin{equation}
\ket{L_{n_-n_+}}
=\frac{1}{\sqrt2}
\Bigl(
\ket{\ell-3}\occ{n_-n_+}
+s_N\ket{3-\ell}\occ{n_+n_-}
\Bigr)\ket{\irr{a}_{1}}_{\mathrm n}.
\label{eq:caoch3_levels}
\end{equation}
The relative sign enforces $\Gmol=\irr{a}_{2}$:
$s_N=-(-1)^{N+k}\parity=(-1)^{N+1}$, using $k=\ell-3$, parity $\parity=(-1)^{v+1}$, and even $\ell+v$.
At $k=0$, $L_{03}$ ($N=3$) and $L_{30}$ ($N=4$) select bend species $\irr{a}_{1}$ and $\irr{a}_{2}$, respectively, realizing Corollary~\ref{cor:shift}. Their subscripts name levels: both use the first term $\ket{0}\occ{03}$ of Eq.~\eqref{eq:caoch3_levels}, with $N=3$ and $N=4$.

The amplitudes in Eq.~\eqref{eq:caoch3_target} describe coherence between complete molecular levels. After tracing out rotation and spin, the bend density operator remains invariant under molecular point group symmetry. Schur's lemma~\cite{hamermesh1989} therefore forbids coherence between inequivalent bend species and requires equal populations and no coherence between the two $\irr{e}$ rows. Only $L_{03}$ and $L_{30}$ individually leave the bend in a pure state.


\begin{table}[h]
\centering
\small
\caption{Preparation-tree levels, all with ortho nuclear spin ($I=\tfrac32$) and $M_F=0$. The $\ell$ and $k$ entries refer to the first component of Eq.~\eqref{eq:caoch3_levels}, except for $L_{03}$ and $L_{30}$, whose components have $k=0$ and $\ell=\pm3$. The state column uses the notation of Table~\ref{tab:ammonia_species}. Energies follow from Eq.~\eqref{eq:caoch3_hamiltonian} alone, with the constants in the upper block of Table~\ref{tab:caoch3_constants}; they include neither centrifugal distortion nor the spin--rotation and hyperfine shifts (Appendix~\ref{app:caoch3_ham}).}
\label{tab:caoch3_levels}
\begin{tabular}{lrrrrrrcccr}
\toprule
Level & $v$ & $\ell$ & $k$ & $N$ & $J$ & $F$ & $(\Grot,\Gspin,\Gvib)$ & state & role & $(E-E_G)/hc$ ($\mathrm{cm^{-1}}$)\\
\midrule
$G=L_{00}$ & 0 & 0 & $-3$ & 3 & $\tfrac72$ & 2 & $(\irr{a}_{2},\irr{a}_{1},\irr{a}_{1})$ & $\catP$ & start & 0.000\\
$L_{01}$ & 1 & 1 & $-2$ & 3 & $\tfrac72$ & 3 & $(\irr{e},\irr{a}_{1},\irr{e})$ & $\catB{1}$ (r:v) & tree & 137.520\\
$L_{02}$ & 2 & 2 & $-1$ & 3 & $\tfrac72$ & 4 & $(\irr{e},\irr{a}_{1},\irr{e})$ & $\catB{1}$ (r:v) & tree & 275.068\\
$L_{11}$ & 2 & 0 & $-3$ & 4 & $\tfrac92$ & 4 & $(\irr{a}_{2},\irr{a}_{1},\irr{a}_{1})$ & $\catP$ & tree & 277.267\\
$L_{03}$ & 3 & $\pm3$ & 0 & 3 & $\tfrac72$ & 5 & $(\irr{a}_{2},\irr{a}_{1},\irr{a}_{1})$ & $\catP$ & target & 412.642\\
$L_{30}$ & 3 & $\pm3$ & 0 & 4 & $\tfrac92$ & 5 & $(\irr{a}_{1},\irr{a}_{1},\irr{a}_{2})$ & $\catP$ & target & 413.573\\
$L_{12}$ & 3 & 1 & $-2$ & 4 & $\tfrac92$ & 5 & $(\irr{e},\irr{a}_{1},\irr{e})$ & $\catB{1}$ (r:v) & target & 412.667\\
$L_{21}$ & 3 & $-1$ & $-4$ & 4 & $\tfrac92$ & 5 & $(\irr{e},\irr{a}_{1},\irr{e})$ & $\catB{1}$ (r:v) & target & 416.110\\
\bottomrule
\end{tabular}
\end{table}

\subsection{Transition dipole matrix elements}
\label{app:caoch3_dipoles}

For an electric field polarized along the laboratory axis $Z$, the transition dipole matrix element $d=\bra{u}\mu_Z\ket{\mathrm{lo}}$ couples the lower and upper levels. We calculate $d$ in four steps: transform $\mu_Z$ to the molecular frame, evaluate the bend and rotational matrix elements, and recouple the spins~\cite{Edmonds1957,Zare1988,BrownCarrington2003}. We follow Ref.~\cite{Edmonds1957} for Wigner functions, angular-momentum coupling coefficients ($3j$ and $6j$ symbols), and reduced matrix elements.

\emph{Frame transformation.} The spherical dipole components are $\mu_{\pm1}=\mp\tfrac{1}{\sqrt2}\,(\mu_x\pm\iunit\mu_y)$ and $\mu_0=\mu_z$~\cite{Zare1988,BrownCarrington2003}. The molecular orientation $\orient$ is the rotation that carries the laboratory axes $(X,Y,Z)$ into the molecule-fixed axes $(x,y,z)$. It is specified by three Euler angles~\cite{Zare1988}: the polar coordinates (polar angle and azimuth) of the molecular $z$ axis in the laboratory frame, and the angle of rotation of the molecule about that axis. Rotating to the laboratory frame gives $\mu_Z=\sum_q [\wignerD{1}_{0q}(\orient)]^*\mu_q$, where $q=0,\pm1$ labels molecular-frame components and $\wignerD{1}_{0q}(\orient)=\bra{1,0}e^{-\iunit\alpha J_z}e^{-\iunit\beta J_y}e^{-\iunit\gamma J_z}\ket{1,q}=e^{-\iunit q\gamma}d^{1}_{0q}(\beta)$ is the Wigner $D$ function of the Euler angles $(\alpha,\beta,\gamma)$ of $\orient$~\cite{Zare1988}. The complex conjugate is needed because $(\alpha,\beta,\gamma)$ orient the molecular axes relative to the laboratory axes. Each coefficient $[\wignerD{1}_{0q}(\orient)]^*$ is then the direction cosine between the laboratory $Z$ axis and the molecular spherical component $q$. In Cartesian form, $\mu_Z=\cos\beta\,\mu_z-\sin\beta\cos\gamma\,\mu_x+\sin\beta\sin\gamma\,\mu_y$, and without the asterisk the $\mu_y$ term would change sign. Thus, a single laboratory polarization includes both $q=+1$ and $q=-1$ components.

\emph{Bend.} Symmetry places the permanent dipole along $z$ and confines its first-order response to an $\irr{e}$ displacement to the $x$--$y$ plane. The dipole derivatives for the $\nu_8$ coordinates $(Q_1,Q_2)$ of Appendix~\ref{app:circular_vibrations} are orthogonal and equal in magnitude~\cite{Herzberg1945,BunkerJensen2006}. Let $\mu_8=\bra{1}\mu_x\ket{0}$ be the fundamental transition dipole moment along one coordinate, and let $a_j$ and $a_j^\dagger$ annihilate and create Cartesian oscillator quanta. The linear dipole operators are then $\mu_x=\mu_8(a_1+a_1^\dagger)$ and $\mu_y=\mu_8(a_2+a_2^\dagger)$. Transforming to the circular operators of Appendix~\ref{app:circular_vibrations} gives $\mu_q=-q\mu_8(a_{-q}+a_q^\dagger)$ for $q=\pm1$.

The creation operator $a_q^\dagger$ adds one bend quantum with vibrational angular momentum $q\hbar$ about $z$, which circulates counterclockwise for $q=+1$ and clockwise for $q=-1$. The annihilation operator $a_{-q}$ removes one quantum circulating in the opposite direction. These processes change $v$ by $+1$ and $-1$, respectively, and, because $\ell=n_+-n_-$, both change $\ell$ by $q$~\cite{Nielsen1951}. Because $\mu_0$ has no first-order bend dependence, $\Delta\ell=0$ bend transitions are absent in this approximation. For an upward transition, the creation-operator matrix element gives
\begin{equation}
V_q=\bra{n_q+1}\mu_q\ket{n_q}=-q\mu_8\sqrt{n_q+1},\qquad n_{+1}=n_+,\quad n_{-1}=n_-.
\label{eq:caoch3_bend_factor}
\end{equation}

\emph{Rotation.} For laboratory projection $M_N$, the normalized symmetric-top wavefunction is $\langle\orient|NkM_N\rangle=[(2N+1)/8\pi^2]^{1/2}\,\wignerD{N}_{M_Nk}(\orient)^*$. Substituting into the matrix element of $\wignerD{1}_{0q}(\orient)^*$ gives an integral of three Wigner functions, yielding two $3j$ symbols~\cite{Edmonds1957,Zare1988}. Primes denote upper-level quantum numbers:
\begin{equation}
\bra{N'k'M_N'}\wignerD{1}_{0q}(\orient)^*\ket{NkM_N}=(-1)^{M_N'-k'}\sqrt{(2N'+1)(2N+1)}\begin{pmatrix}N'&1&N\\-M_N'&0&M_N\end{pmatrix}\begin{pmatrix}N'&1&N\\-k'&q&k\end{pmatrix}.
\label{eq:caoch3_rotation_me}
\end{equation}
Both symbols require $N'=N$ or $N\pm1$. The first gives the laboratory-projection dependence, and the second requires $k'=k+q$. Together with the bend selection rule, this gives the perpendicular-band rule $\Delta k=\Delta\ell=q$~\cite{Herzberg1945}, which conserves $k-\ell$ exactly. Molecular point group symmetry requires only conservation of $k-\ell\bmod3$, the residue determining compatibility with ortho or para spin (Appendix~\ref{subsec:tour}), because the dipole preserves nuclear spin.

\emph{Spins.} Neither $\mu_q$ nor $\wignerD{1}_{0q}$ acts on electron spin $S$ or proton spin $I$, so the ortho state $\ket{\irr{a}_{1}}_{\mathrm{n}}$ is conserved. Conservation of $k-\ell$ restricts the branch and mirror components of Eq.~\eqref{eq:caoch3_levels} to their respective upper-level partners.

In the coupled basis, an operator acting on rotation acquires one $6j$ symbol for each spin recoupling~\cite{Edmonds1957,BrownCarrington2003}. The two recouplings, $\mathbf J=\mathbf N+\mathbf S$ and $\mathbf F=\mathbf J+\mathbf I$, impose $\Delta J=0,\pm1$ and $\Delta F=0,\pm1$. The Wigner--Eckart theorem~\cite{Edmonds1957} supplies the dependence on the laboratory projection $M_F$ and requires $\Delta M_F=0$ for $Z$ polarization.

Combining these steps gives the branch-component matrix element $d_a$, before including the normalization of Eq.~\eqref{eq:caoch3_levels}:
\begin{align}
d_a={}&f_Mf_If_Sf_N\,V_q,\qquad
f_M=(-1)^{F'-M_F}\begin{pmatrix}F'&1&F\\-M_F&0&M_F\end{pmatrix},\nonumber\\
f_I={}&(-1)^{J'+I+F+1}\sqrt{(2F'+1)(2F+1)}\begin{Bmatrix}J'&F'&I\\F&J&1\end{Bmatrix},\qquad
f_S=(-1)^{N'+S+J+1}\sqrt{(2J'+1)(2J+1)}\begin{Bmatrix}N'&J'&S\\J&N&1\end{Bmatrix},\nonumber\\
f_N={}&(-1)^{N'-k'}\sqrt{(2N'+1)(2N+1)}\begin{pmatrix}N'&1&N\\-k'&q&k\end{pmatrix},
\label{eq:caoch3_dipole_factors}
\end{align}
where the curly (round) brackets denote Wigner $6j$ ($3j$) symbols~\cite{Edmonds1957}, and $f_N$ is the reduced rotational matrix element from Eq.~\eqref{eq:caoch3_rotation_me}. The factors $f_S$ and $f_I$ account for spin recoupling, and $f_M$ restores the laboratory-projection dependence.



\subsection{Pulse sequence and experimental requirements}
\label{app:caoch3_pulses}

Using the transition frequencies and dipole matrix elements, we construct selective pulses that distribute population through the tree and set the target amplitudes' relative phases.

\emph{Gate.} For pulse area $\Theta$ and phase $\phi$, the gate in Eq.~\eqref{eq:caoch3_gate} has the following matrix in the ordered basis $(\ket{\mathrm{lo}},\ket{u})$:
\begin{equation}
R_{u,\mathrm{lo}}(\Theta,\phi)=\begin{pmatrix}\cos\tfrac\Theta2&-\iunit\expe^{-\iunit\phi}\sin\tfrac\Theta2\\[2pt]-\iunit\expe^{\iunit\phi}\sin\tfrac\Theta2&\cos\tfrac\Theta2\end{pmatrix},
\label{eq:caoch3_gate_matrix}
\end{equation}
The gate leaves all other levels unchanged and transfers amplitude in both directions along the driven link. On the Bloch sphere, which represents two-level pure states, it rotates by $\Theta$ about an equatorial axis of azimuth $\phi$~\cite{AllenEberly1975,Nielsen_Chuang_2010}. A $\pi$ pulse exchanges the levels with phases $\ket{\mathrm{lo}}\to-\iunit\expe^{\iunit\phi}\ket{u}$ and $\ket{u}\to-\iunit\expe^{-\iunit\phi}\ket{\mathrm{lo}}$. A $\pi/2$ pulse creates an equal-weight superposition from either level. Since $R_{u,\mathrm{lo}}(\Theta,\phi)^{-1}=R_{u,\mathrm{lo}}(\Theta,\phi+\pi)$, reversing the pulse order and shifting each phase by $\pi$ inverts the sequence.

\emph{Physical realization.} A resonant field $\mathcal E(t)\cos(\omega t-\phi_L)$ implements this gate in the rotating-wave approximation~\cite{AllenEberly1975}, which neglects rapidly oscillating, nonresonant terms. Here $\mathcal E(t)$ is the field envelope, $\omega$ the angular frequency, and $\phi_L$ the field phase. The Rabi frequency $\Omega(t)=|d|\mathcal E(t)/\hbar$ sets the pulse area $\Theta=\int\Omega\,\diff t$, while $\phi=\phi_L+\arg(-d)$. Thus, $\phi_L=\phi$ for $d<0$ and $\phi_L=\phi+\pi$ for $d>0$. Only P5 has negative $d$ (Table~\ref{tab:caoch3_pulses}).

\emph{Requirements.} Realistic implementation of Eq.~\eqref{eq:caoch3_gate_matrix} requires five experimental conditions.

(i) Resonance: the tone must match the transition angular frequency $\omega_{u,\mathrm{lo}}$ with a detuning much smaller than $\Omega$.

(ii) Selectivity: every competing transition from a populated level must have angular-frequency detuning $\Delta\gg\Omega,\tau^{-1}$, where $\tau$ is the pulse duration. Off-resonant population errors are then of order $(\Omega/\Delta)^2$~\cite{AllenEberly1975}. Excluding the line that coincides with P4, which condition (iv) avoids, the nearest competing rovibrational line lies $\Delta/2\pi=0.73$~GHz from P4 (Appendix~\ref{app:caoch3_ham}). At the Rabi frequencies chosen below, all at most $44$~kHz, the estimated population error is below $4\times10^{-9}$.

Spin structure imposes a tighter constraint. It is the only term beyond Eq.~\eqref{eq:caoch3_hamiltonian} that sets a parameter of the protocol. We estimate it using the measured ground-state spin--rotation and proton hyperfine constants~\cite{Namiki1998} for all bend levels (lower block of Table~\ref{tab:caoch3_constants}). Each rovibrational line splits into components, including unwanted transitions from populated levels. The nearest lies $0.36$~MHz from P7 and has comparable strength. Requiring $\Omega/2\pi\le12$~kHz keeps its estimated population error below $10^{-3}$. The same estimate bounds the other six pulses less tightly. We therefore drive all seven links at one intensity, chosen so that P7 has $\Omega/2\pi=10$~kHz. Because $\Omega\propto|d|$, the other links then have $\Omega/2\pi$ between $9.8$ and $44$~kHz, each below its own bound, and every estimated population error stays below $7\times10^{-4}$.

(iii) Phase reference: all tones must share a clock defining the rotating frame of the target amplitudes.

(iv) Order: at each branch point, the partial pulse must precede the $\pi$ pulse, and each parent must be populated before driving its children. P4 must finish before P5 begins, so that $L_{30}$ stays empty throughout P4 and the line that coincides with P4 (Appendix~\ref{app:caoch3_ham}) has no population to drive. Gates on disjoint level pairs commute, allowing P6 and P7 to overlap with P4 and P5.

(v) Coherence: the sequence duration, at most $0.22$~ms here (see below), must be short compared with the bend-level lifetimes.

\emph{Intensity and duration.} Converting these Rabi frequencies into a laser intensity requires the absolute bend transition dipole $\mu_8$. The linear dipole operator $\mu_x=\mu_8(a_1+a_1^\dagger)$ of Appendix~\ref{app:caoch3_dipoles} is the first-order term $(\partial\mu_x/\partial Q_1)\,Q_1$, where the mass-weighted normal coordinate is $Q_1=(\hbar/2\omega)^{1/2}(a_1+a_1^\dagger)$ and $\omega$ is the harmonic angular frequency of $\nu_8$. Hence $\mu_8=(\hbar/2\omega)^{1/2}\,\partial\mu_x/\partial Q_1$~\cite{Herzberg1945,BunkerJensen2006}. We evaluate the derivative by central differences of the dipole moment along both $\nu_8$ components of the B3LYP force field of Appendix~\ref{app:caoch3_ham}, whose harmonic wavenumber is $150.7\ \mathrm{cm^{-1}}$, and obtain $\mu_8=0.150$~D. The same calculation gives a permanent dipole of $1.45$~D, $8\%$ below the measured $1.58(8)$~D~\cite{Namiki1998}. P7, with $|d|=0.124\,\mu_8$, then reaches $\Omega/2\pi=10$~kHz at the field amplitude $\mathcal E=107$~V/m, that is, at the intensity $I=c\varepsilon_0\mathcal E^2/2=1.5\ \mathrm{mW/cm^2}$~\cite{Foot2005}. We assume this intensity for all seven tones. A pulse of area $\Theta$ on link P$n$ then lasts $\Theta/\Omega_n$. The last column of Table~\ref{tab:caoch3_pulses} lists the $\pi$-pulse durations $\tau_\pi=\pi/\Omega_n$, from $11$ to $51~\mu\mathrm{s}$. With the pulses applied one after another, the single-SALC sequences (b) and (c) last $91$ to $112~\mu\mathrm{s}$. The equal superposition with every $c=\tfrac12$ takes $171~\mu\mathrm{s}$, or $145~\mu\mathrm{s}$ when P6 and P7 overlap P4 and P5 [condition (iv)]. No target in (a) takes longer than $217~\mu\mathrm{s}$, the sum of all seven $\tau_\pi$.

\emph{Solution.} We determine the pulse parameters in Eq.~\eqref{eq:caoch3_sequence} by working backward from the target amplitudes to $G$. At each branch, a partial pulse transfers amplitude from a parent $\mathrm{lo}$ to an initially empty child $u$. A subsequent $\pi$ pulse transfers the remainder to the other initially empty child $u'$. For initial parent amplitude $\beta$, Eq.~\eqref{eq:caoch3_gate_matrix} gives
\begin{equation}
R_{u',\mathrm{lo}}(\pi,\phi')\,R_{u,\mathrm{lo}}(\Theta,\phi)\,\beta\ket{\mathrm{lo}}=-\iunit\beta\bigl(\expe^{\iunit\phi}\sin\tfrac\Theta2\ket{u}+\expe^{\iunit\phi'}\cos\tfrac\Theta2\ket{u'}\bigr).
\label{eq:caoch3_split}
\end{equation}
For desired child amplitudes $\alpha$ and $\alpha'$, population conservation requires $|\beta|^2=|\alpha|^2+|\alpha'|^2$. Matching magnitudes gives $\tan\tfrac\Theta2=|\alpha|/|\alpha'|$, and matching phases gives $\phi=\arg\alpha-\arg\beta+\tfrac\pi2$ and $\phi'=\arg\alpha'-\arg\beta+\tfrac\pi2$. A parent with one child requires only a $\pi$ pulse. Both cases empty the parent, leaving no population in intermediate levels.

Write each target amplitude in Eq.~\eqref{eq:caoch3_target} as $c=|c|\,\expe^{\iunit\gamma}$, with $\gamma$ carrying the labels of $c$. The branch populations are $W_{\irr{a}}=|c^{\irr{a}_{1}}|^2+|c^{\irr{a}_{2}}|^2$ and $W_{\irr{e}}=|c^{\irr{e}}_{1}|^2+|c^{\irr{e}}_{2}|^2$. Applying Eq.~\eqref{eq:caoch3_split} gives the partial-pulse areas for P2, P4, and P6 in Eq.~\eqref{eq:caoch3_areas}, starting with $\tan\tfrac{\Theta_2}{2}=(W_{\irr{a}}/W_{\irr{e}})^{1/2}$. P1, P3, P5, and P7 are $\pi$ pulses. With $\phi_1=\phi_2=\phi_3=0$, successive applications give
\begin{align}
R_1\ket{G}&=-\iunit\ket{L_{01}},\nonumber\\
R_3R_2R_1\ket{G}&=-\bigl(\sqrt{W_{\irr{a}}}\ket{L_{02}}+\sqrt{W_{\irr{e}}}\ket{L_{11}}\bigr),\nonumber\\
R_7\cdots R_1\ket{G}&=\iunit\bigl(|c^{\irr{a}_{1}}|\,\expe^{\iunit\phi_4}\ket{L_{03}}+|c^{\irr{a}_{2}}|\,\expe^{\iunit\phi_5}\ket{L_{30}}+|c^{\irr{e}}_{1}|\,\expe^{\iunit\phi_7}\ket{L_{21}}+|c^{\irr{e}}_{2}|\,\expe^{\iunit\phi_6}\ket{L_{12}}\bigr).
\label{eq:caoch3_layers}
\end{align}
The final line equals $\iunit\ket{\Psi_T}$ when $\phi_4=\gamma^{\irr{a}_{1}}$, $\phi_5=\gamma^{\irr{a}_{2}}$, $\phi_6=\gamma^{\irr{e}}_{2}$, and $\phi_7=\gamma^{\irr{e}}_{1}$. Three adjustable areas and three relative phases supply the six real parameters of a normalized four-level target, excluding its global phase. The sequence therefore reaches every superposition in the target subspace. 

\begin{table}[h]
\centering
\small
\caption{Pulse sequences for Fig.~\ref{fig:caoch3_tree}, with gate area $\Theta_n$ and phase $\phi_n$ for $R_n$ [Eq.~\eqref{eq:caoch3_gate}]; a dash denotes an omitted pulse. (a) Arbitrary target, using the areas in Eq.~\eqref{eq:caoch3_areas} and target phases $\gamma$. (b) $L_{30}$, labeled by $\ket{\irr{a}_{2}}$; replacing P5 with P4 prepares $L_{03}$, labeled by $\ket{\irr{a}_{1}}$. (c) $L_{21}$, labeled by $\irket{\irr{e}}{1}$; replacing P7 with P6 prepares $L_{12}$, labeled by $\irket{\irr{e}}{2}$. Each sequence yields $\iunit\ket{\Psi_T}$. Couplings follow Eq.~\eqref{eq:caoch3_dipole_factors}; line positions follow Eq.~\eqref{eq:caoch3_lines} with the constants in the upper block of Table~\ref{tab:caoch3_constants}. They are rounded after the calculation, so they can differ by $0.001\ \mathrm{cm^{-1}}$ from differences of the rounded energies in Table~\ref{tab:caoch3_levels}. The last column gives the $\pi$-pulse duration $\tau_\pi=\pi/\Omega$ at the common intensity $I=1.5\ \mathrm{mW/cm^2}$ assumed in Appendix~\ref{app:caoch3_pulses}; a pulse of area $\Theta$ lasts $(\Theta/\pi)\,\tau_\pi$.}
\label{tab:caoch3_pulses}
\begin{tabular}{llrrrrcccr}
\toprule
Pulse & Link & $q$ & $\Delta N$ & $d/\mu_8$ & $\tilde\nu$ ($\mathrm{cm^{-1}}$) & (a) arbitrary & (b) $\ket{\irr{a}_{2}}$ & (c) $\irket{\irr{e}}{1}$ & $\tau_\pi$ ($\mu\mathrm{s}$)\\
\midrule
P1 & $G\to L_{01}$ & $+1$ & $0$ & $+0.121$ & 137.520 & $\pi,\ 0$ & $\pi,\ 0$ & $\pi,\ 0$ & 50.9\\
P2 & $L_{01}\to L_{02}$ & $+1$ & $0$ & $+0.220$ & 137.547 & $\Theta_2,\ 0$ & $\pi,\ 0$ & -- & 28.1\\
P3 & $L_{01}\to L_{11}$ & $-1$ & $+1$ & $+0.545$ & 139.746 & $\pi,\ 0$ & -- & $\pi,\ 0$ & 11.4\\
P4 & $L_{02}\to L_{03}$ & $+1$ & $0$ & $+0.231$ & 137.575 & $\Theta_4,\ \gamma^{\irr{a}_{1}}$ & -- & -- & 26.7\\
P5 & $L_{02}\to L_{30}$ & $+1$ & $+1$ & $-0.509$ & 138.505 & $\pi,\ \gamma^{\irr{a}_{2}}$ & $\pi,\ 0$ & -- & 12.2\\
P6 & $L_{11}\to L_{12}$ & $+1$ & $0$ & $+0.164$ & 135.400 & $\Theta_6,\ \gamma^{\irr{e}}_{2}$ & -- & -- & 37.8\\
P7 & $L_{11}\to L_{21}$ & $-1$ & $0$ & $+0.124$ & 138.843 & $\pi,\ \gamma^{\irr{e}}_{1}$ & -- & $\pi,\ 0$ & 50.0\\
\bottomrule
\end{tabular}
\end{table}

\emph{Readout.} Each SALC label identifies a distinct molecular target level. The weights $|c^{\irr{a}_{1}}|^2$, $|c^{\irr{a}_{2}}|^2$, $|c^{\irr{e}}_{1}|^2$, and $|c^{\irr{e}}_{2}|^2$ are therefore molecular populations, distinguishable by state-selective detection because the levels differ in $N$ or $k$.

Relative phases are measured by reversing part of the tree to interfere target amplitudes, as in Ramsey interferometry. Apply P5 as a $\pi$ pulse of phase $\phi_5$, then P4 as a $\pi/2$ pulse of phase $\phi$. The resulting population in $L_{02}$ is $P_{02}=\tfrac12|c^{\irr{a}_{1}}+\expe^{\iunit\varphi}c^{\irr{a}_{2}}|^2$, where $\varphi=\phi-\phi_5$ and phase information can be extracted.  These measurements probe molecular coherences, which disappear when rotation and spin are traced out and cannot be detected by measuring the bend alone.

